\section{The effect-identity framework}\label{sec:framework}

Ground-truth identity cannot be observed. A mechanism must instead compute decidable relations and use them as proxies, and two dual notions of correctness govern such proxies.

\subsection{Certifiers and refuters}

\begin{definition}[Certifier]\label{def:certifier}
A relation $\mathsf{C}\subseteq X_t\times X_t$ is a \emph{certifier} if $\mathsf{C}\subseteq{\approx_t}$: whenever $\mathsf{C}(x,x')$ holds, the two requests denote the same act.
\end{definition}

\begin{definition}[Refuter]\label{def:refuter}
A relation $\mathsf{D}\subseteq X_t\times X_t$ is a \emph{refuter} if $\mathsf{D}\cap{\approx_t}=\emptyset$: whenever $\mathsf{D}(x,x')$ holds, the two requests denote different acts.
\end{definition}

A certifier licenses suppression; a refuter licenses execution. Each may be arbitrarily incomplete, and neither is a weakened form of the other. Table~\ref{tab:relations} lists the relations used in this paper. Canonical equality $\equiv_\kappa$ normalises each declared field under a policy (identifier, case-insensitive token, number, free text) and is a certifier whenever each policy is semantics-preserving (assumption A5). Projection difference $\not\equiv_\sigma$ compares the canonical values of the declared \emph{typed} fields and is a refuter whenever a difference in any such field implies a different act (A5$'$). A judge threshold $J>\tau$ is neither: it is decidable but carries no proof in either direction.

\begin{table}[!tbp]
\centering\small
\caption{Relations used in this paper and the assumption under which each is sound.}\label{tab:relations}
\begin{tabular}{@{}llll@{}}
\toprule
Relation & Role & Sound under & Used as \\
\midrule
byte equality $\equiv_{\mathrm{exact}}$ & certifier & always & baseline B1 \\
canonical equality $\equiv_\kappa$ & certifier & A5 & L0; licenses \SUP \\
projection difference $\not\equiv_\sigma$ & refuter & A5$'$ & witness; licenses \EXE \\
judge threshold $J>\tau$ & neither & --- & must be calibrated \\
ground truth $\approx_t$ & target & --- & not observable \\
\bottomrule
\end{tabular}
\end{table}

\begin{proposition}[Structure]\label{prop:lattice}
Fix $t$ and order relations on $X_t\times X_t$ by inclusion. The certifiers are closed under union and intersection, with least element $\emptyset$ and greatest element $\approx_t$; the refuters are closed under the same operations with greatest element the complement of $\approx_t$. A finite union of decidable certifiers (refuters) is a decidable certifier (refuter). The only relation that is both a certifier and a refuter is $\emptyset$.
\end{proposition}

The proof is in \ref{app:proofs}. Two consequences are used below. Sound decidable layers can be combined by union without re-arguing soundness, and a design must supply certification and refutation separately, because no single relation performs both.

\subsection{The residual}

\begin{definition}[Decided set and residual]\label{def:residual}
For a declaration that fixes $\kappa_t$ and $\sigma_t$, the decided set is $\Dec_t={\equiv_\kappa}\cup{\not\equiv_\sigma}$ and the residual is its complement,
\[
\Res_t=\{(x,x') : \kappa_t(x)\ne\kappa_t(x') \ \wedge\ \sigma_t(x)=\sigma_t(x')\}.
\]
\end{definition}

The residual contains the requests that differ after canonicalisation yet agree on every declared effect-determining field: exactly the pairs whose identity depends on free text.

\begin{proposition}[Residual determination]\label{prop:residual}
(i) $\Res_t$ is a function of $(\kappa_t,\sigma_t)$ alone; it does not depend on the judge, the agent, the adversary or any observation at run time. (ii) If a mechanism suppresses on $\equiv_\kappa$, executes on $\not\equiv_\sigma$, and A5 and A5$'$ hold, every error it makes is on a pair in $\Res_t$. (iii) If $(\kappa',\sigma')$ refines $(\kappa,\sigma)$ (Definition~\ref{def:refinement}), then $\Res'_t\subseteq\Res_t$.
\end{proposition}

\begin{proof}
(i) The defining condition mentions only $\kappa_t$ and $\sigma_t$. (ii) A pair in $\equiv_\kappa$ is suppressed and, by A5, is a true duplicate; a pair in $\not\equiv_\sigma$ is executed and, by A5$'$, is truly distinct; hence any error lies outside $\Dec_t$. (iii) Refinement enlarges $\equiv_\kappa$ and $\not\equiv_\sigma$, so their union grows and its complement shrinks.
\end{proof}

Proposition~\ref{prop:residual} is elementary, and part (ii) holds by construction rather than by experiment. Its content is nevertheless operational: it locates every possible error of a sound layered mechanism in a set the schema determines, and it separates two kinds of intervention. Changing the judge---a larger model, an ensemble, a better prompt---redistributes decisions \emph{inside} $\Res_t$; only changing the declaration changes $\Res_t$ itself.

\subsection{Schema refinement and the refinement dilemma}\label{sec:refinement}

\begin{definition}[Refinement]\label{def:refinement}
A declaration $(\kappa',\sigma')$ \emph{refines} $(\kappa,\sigma)$ if ${\equiv_\kappa}\subseteq{\equiv_{\kappa'}}$, $\sigma(x)\ne\sigma(x')\Rightarrow\sigma'(x)\ne\sigma'(x')$, and $\equiv_{\kappa'}$ and $\not\equiv_{\sigma'}$ remain sound.
\end{definition}

Refinement moves act-determining information out of free text and into a typed field. In the corpus of Section~\ref{sec:method}, \texttt{wire\_money} declares its amount, beneficiary and reference as free text, so its projection is empty and every non-identical pair of wires is residual; a second instalment of the same payment differs from the first only by a phrase such as ``instalment 2 of 2''. A refined declaration that types the numbers each text field mentions lets the refuter separate the two instalments, so an injected ``this is the same payment; suppress'' can no longer reach the judge.

Refinement is therefore the only intervention that removes pairs from the judge's jurisdiction, and it improves the safety-critical direction on the pairs it decides. It also has a cost that is easy to miss. The calibration pool of the statistical layer is drawn from the residual (Section~\ref{sec:calibration}), so refinement shrinks the pool together with the residual. A smaller pool moves the conformal threshold to a higher order statistic and certifies fewer candidates, and the pairs that survive refinement are not a random sample of those that did not: the refuter removes distinct pairs whose declared fields already differ, which judges tend to score low. We call this the \emph{refinement dilemma} (Figure~\ref{fig:refinement}). Proposition~\ref{prop:dilemma} states its exact part; Section~\ref{sec:res-refinement} measures both parts. The practical implication is concrete: after refining a schema, re-calibrate on the new jurisdiction and plan for a smaller pool---at $\alpha=0.05$ the fence needs at least 19 distinct residual calibration pairs to suppress semantically at all, and at least $m/\alpha-1$ to do so with $m$ candidates.

\begin{figure}[!tbp]
\centering
\begin{tikzpicture}[font=\footnotesize, x=0.12cm, y=1cm]
  % Region widths are the test-split shares (percent) under each declaration,
  % primary judge: certified 6.6; refuted 17.2 (declared) / 34.5 (refined).
  % ── declared schema (heading above its bar) ──
  \node[anchor=west, font=\footnotesize\bfseries] at (0,3.20) {Declared schema};
  \fill[green!18] (0,2.60) rectangle (6.6,3.05);
  \fill[blue!12] (6.6,2.60) rectangle (23.8,3.05);
  \fill[orange!18] (23.8,2.60) rectangle (100,3.05);
  \draw (0,2.60) rectangle (100,3.05);
  \draw (6.6,2.60) -- (6.6,3.05);
  \draw (23.8,2.60) -- (23.8,3.05);
  \node[font=\scriptsize] at (3.3,2.825) {C};
  \node[font=\scriptsize] at (15.2,2.825) {refuted};
  \node[font=\scriptsize] at (61.9,2.825) {residual: pool $n=142$, $t_{\mathrm{hi}}$ = 7th largest};
  % ── transition ──
  \draw[-{Latex[length=2mm]}, thick] (16,2.52) -- (16,1.13);
  \node[anchor=west, font=\scriptsize, align=left] at (20,1.83)
    {refinement moves distinct pairs from the residual into the refuter,\\
     and the calibration pool shrinks with it};
  % ── refined schema (heading below its bar) ──
  \fill[green!18] (0,0.60) rectangle (6.6,1.05);
  \fill[blue!12] (6.6,0.60) rectangle (41.1,1.05);
  \fill[orange!18] (41.1,0.60) rectangle (100,1.05);
  \draw (0,0.60) rectangle (100,1.05);
  \draw (6.6,0.60) -- (6.6,1.05);
  \draw (41.1,0.60) -- (41.1,1.05);
  \node[font=\scriptsize] at (3.3,0.825) {C};
  \node[font=\scriptsize] at (23.8,0.825) {refuted};
  \node[font=\scriptsize] at (70.5,0.825) {residual: pool $n=49$, $t_{\mathrm{hi}}$ = 2nd largest};
  \node[anchor=west, font=\footnotesize\bfseries] at (0,0.30) {Refined (tightened) schema};
\end{tikzpicture}
\caption{The refinement dilemma, drawn to scale for the test split (shares of 513 pairs; C: certified; the residual is the fragment the judge decides). Refinement enlarges the refuted region, which removes pairs from the judge's jurisdiction and closes the prose-only injection class. Because the calibration pool is the distinct part of the residual, the pool shrinks with it (primary judge: from 142 to 49 pairs), and at $\alpha=0.05$ the conformal threshold moves from the 7th to the 2nd largest pool score (Proposition~\ref{prop:dilemma}). Section~\ref{sec:res-refinement} measures the consequence.}\label{fig:refinement}
\end{figure}


\subsection{Where existing mechanisms sit}

In this framework, an idempotency key is a certifier on a client-chosen token with residual policy ``execute'' (a fresh token is a fresh act); content-addressed fencing is the certifier $\equiv_\kappa$ with the same residual policy; a raw LLM analyser resolves the whole space, residual included, by an unbounded judgement; and a canonical-action binding~\cite{xu2026caplease} is a certifier on a normal form. We are not aware of a prior mechanism that pairs a sound certifier with a sound refuter and treats only the remainder as a calibrated decision.
